\section{Lab 5: AirSim Drone Simulation and Navigation Control}

\subsection{Objective}
In this lab, we worked with AirSim, a powerful drone simulator, to learn how to:
\begin{itemize}
	\item Set up and configure a virtual drone environment
	\item Control a simulated drone using our keyboard
	\item Monitor the drone's position, speed, and orientation in real-time
	\item Record detailed flight data for later analysis
	\item Create visual plots to understand how the drone behaved during flight
	\item Successfully navigate the drone from a starting point to a target location
\end{itemize}

\subsection{Requirements}
To complete this lab, we needed:
\begin{itemize}
	\item \textbf{Operating System:} Windows or Linux (compatible with AirSim)
	\item \textbf{Simulator:} AirSim running on Unreal Engine or Unity
	\item \textbf{Programming:} Python 3.x with AirSim API
	\item \textbf{Python Libraries:} \texttt{airsim}, \texttt{keyboard}, \texttt{pandas}, \texttt{matplotlib}, \texttt{numpy}
	\item \textbf{Computer:} A system powerful enough to run 3D simulations smoothly
\end{itemize}

\subsection{Theory}

\subsubsection{What is AirSim?}

AirSim is an open-source simulator built on Unreal Engine that lets us fly virtual drones without the risk of crashing expensive hardware. Think of it as a high-fidelity flight simulator for autonomous vehicles. It offers some really useful features:
\begin{itemize}
	\item Realistic physics that mimics how actual drones fly
	\item Multiple virtual sensors including cameras, GPS, and IMUs
	\item Python and C++ APIs so we can write code to control the drone
	\item Support for testing with real flight controllers
	\item Customizable environments and different types of vehicles
\end{itemize}

\subsubsection{Understanding the Coordinate System}

AirSim uses what's called the NED (North-East-Down) coordinate system. At first, this seems a bit backwards, but there's a reason for it:
\begin{itemize}
	\item \textbf{X-axis:} Points North (the drone's forward direction)
	\item \textbf{Y-axis:} Points East (to the drone's right)
	\item \textbf{Z-axis:} Points Down (yes, down is positive!)
\end{itemize}

The tricky part here is the Z-axis. When the drone goes up, the Z value becomes more negative. When it descends, Z becomes more positive. It takes a bit of getting used to, but it's a standard convention in aviation and robotics.

\subsubsection{How We Control the Drone}

AirSim gives us several ways to control our virtual drone:
\begin{itemize}
	\item \textbf{Velocity Control:} We can tell the drone to move at specific speeds in different directions
	\item \textbf{Position Control:} We can command the drone to fly to exact coordinates
	\item \textbf{Angle Control:} We can directly control the drone's tilt, rotation, and throttle
	\item \textbf{RC Control:} We can simulate using a remote controller
\end{itemize}

For our lab, we chose velocity control combined with yaw (rotation) control because it feels the most natural when flying with a keyboard—kind of like a first-person video game.

\subsubsection{What Data Can We Collect?}

The simulated drone has access to various sensors, just like a real drone:
\begin{itemize}
	\item \textbf{Kinematics:} Where the drone is, how fast it's moving, and which way it's facing
	\item \textbf{GPS:} Geographic coordinates (latitude, longitude, altitude)
	\item \textbf{IMU:} Acceleration and rotation rates
	\item \textbf{Barometer:} Altitude based on air pressure
	\item \textbf{Magnetometer:} Compass heading
\end{itemize}

\subsection{Procedure}

\subsubsection{Setting Up AirSim}

The first step was creating a configuration file called \texttt{settings.json}. This file tells AirSim how to set up our simulation. We placed it in the AirSim folder:
\begin{itemize}
	\item \textbf{On Windows:} \texttt{C:/Users/[Username]/Documents/AirSim/}
	\item \textbf{On Linux:} \texttt{\textasciitilde/Documents/AirSim/}
\end{itemize}

\textbf{Our configuration file (settings.json):}
\begin{lstlisting}[language=JSON]
{
  "SeeDocsAt": "https://github.com/Microsoft/AirSim/blob/main/docs/settings.md",
  "SettingsVersion": 1.2,
  "SimMode": "Multirotor",
  "Vehicles": {
    "Drone1":{
      "VehicleType": "SimpleFlight",
      "AutoCreate": true,
      "Cameras":{
        "FrontCam": {
          "CaptureSettings":[
            {"ImageType":0, "Width": 1280, "Height":720, "FOV_Degrees":90},
            {"ImageType":5, "Width": 1280, "Height":720, "FOV_Degrees":90}
          ],
          "X": 2.0, "Y": 0.0, "Z": -0.2,
          "Pitch": 0.0, "Roll": 0.0, "Yaw": 0.0
        }
      },
      "X": 0.0, "Y": 0.0, "Z": -0.5, "Yaw": 0.0
    }
  },
  "Wind": {"X": 3.0, "Y": 0.0, "Z": 0.0}
}
\end{lstlisting}

\textbf{What these settings mean:}
\begin{itemize}
	\item \texttt{"SimMode": "Multirotor"}: We're simulating a quadcopter-style drone
	\item \texttt{"VehicleType": "SimpleFlight"}: We're using AirSim's built-in, easy-to-use flight controller
	\item \texttt{"AutoCreate": true}: The drone appears automatically when we start the simulation
	\item \textbf{Camera setup:}
	\begin{itemize}
		\item We added a front-facing camera that captures HD video (1280x720)
		\item It has a 90-degree field of view, similar to an action camera
		\item The camera is mounted 2 meters in front of the drone, slightly above center
		\item We set it up to capture both regular RGB images and segmentation data for computer vision
	\end{itemize}
	\item \textbf{Starting position:}
	\begin{itemize}
		\item The drone spawns at the origin (0, 0) coordinates
		\item It starts 0.5 meters above the ground (Z = -0.5, remember down is positive!)
		\item It faces North (yaw = 0 degrees)
	\end{itemize}
	\item \textbf{Wind conditions:}
	\begin{itemize}
		\item We added a 3 m/s wind blowing from the North
		\item This makes the simulation more realistic—the drone has to compensate for drift
	\end{itemize}
\end{itemize}

\subsubsection{Building the Drone Control Script}

Next, we wrote a Python script that lets us fly the drone using our keyboard. Here's how we built it, piece by piece:

\textbf{1. Connecting to the Simulator:}
\begin{lstlisting}[language=Python]
import airsim
import keyboard
import time 
import csv
import math 
import os
from datetime import datetime

client = airsim.MultirotorClient()
client.confirmConnection()
\end{lstlisting}

This code imports the libraries we need and establishes a connection to the AirSim simulator. We also imported the \texttt{keyboard} library so we can detect key presses, and \texttt{csv} for saving our flight data.

\textbf{2. Setting Up Flight Data Recording:}
\begin{lstlisting}[language=Python]
os.makedirs("flightLogs", exist_ok=True)
log_path = f"flightLogs/drone_path_{datetime.now().strftime('%Y%m%d_%H%M%S')}.csv"
logfile = open(log_path, "w", newline="")
writer = csv.writer(logfile)

writer.writerow([
    "time", "command", "vx", "vy", "vz", "yawRate",
    "x", "y", "z", "roll", "pitch", "yaw",
    "gpsLat", "gpsLon", "gpsAlt"
])
\end{lstlisting}

We created a folder to store our flight logs and set up a CSV file with a timestamp in its name. Each log file records 15 different parameters every 0.05 seconds (20 times per second), giving us a detailed record of everything the drone did.

\textbf{3. Defining Control Parameters:}
\begin{lstlisting}[language=Python]
Vel = 2.5          # How fast the drone moves (meters/second)
Yaw_Rate = 25      # How fast the drone rotates (degrees/second)
vx = vy = vz = 0.0 # Current velocity commands
YawRate = 0.0      # Current rotation rate
currentCmd = "idle"
flying = False
running = True
\end{lstlisting}

We chose a cruising speed of 2.5 m/s, which gives us smooth, controllable flight. The rotation rate of 25 degrees per second lets us turn quickly without losing control.

\textbf{4. The Main Control Loop:}
\begin{lstlisting}[language=Python]
while running:
    vx = vy = vz = 0.0
    YawRate = 0.0
    currentCmd = "idle"
    time.sleep(0.05)  # Update 20 times per second

    # Takeoff sequence
    if keyboard.is_pressed("y") and not flying:
        client.enableApiControl(True)
        client.armDisarm(True)
        client.takeoffAsync().join()
        client.moveToZAsync(-10, 2).join()
        flying = True
        currentCmd = "takeoff"
    
    # Landing sequence
    if keyboard.is_pressed("l") and flying:
        client.landAsync().join()
        client.armDisarm(False)
        client.enableApiControl(False)
        flying = False
        currentCmd = "land"

    if not flying:
        time.sleep(0.05)
        continue
\end{lstlisting}

\textbf{How takeoff works:}
\begin{itemize}
	\item We press 'Y' to start the sequence
	\item First, we enable API control (giving our code permission to fly the drone)
	\item Then we arm the motors (like turning the key in a car)
	\item The drone automatically lifts off to a safe altitude
	\item Finally, it climbs to 10 meters (Z = -10, remember the coordinate system!)
\end{itemize}

\textbf{5. Keyboard Controls for Movement:}
\begin{lstlisting}[language=Python]
    # Forward and backward
    if keyboard.is_pressed("w"):
        vx = Vel
        currentCmd = "forward"
    elif keyboard.is_pressed("s"):
        vx = -Vel
        currentCmd = "backward"
    
    # Left and right
    if keyboard.is_pressed("a"):
        vy = -Vel
        currentCmd = "left"
    elif keyboard.is_pressed("d"):
        vy = Vel
        currentCmd = "right"
    
    # Up and down
    if keyboard.is_pressed("z"):
        vz = -Vel
        currentCmd = "up"
    elif keyboard.is_pressed("c"):
        vz = Vel
        currentCmd = "down"
    
    # Rotation controls
    if keyboard.is_pressed("q"):
        YawRate = Yaw_Rate
        currentCmd = "yawLeft"
    elif keyboard.is_pressed("e"):
        YawRate = -Yaw_Rate
        currentCmd = "yawRight"

    if keyboard.is_pressed("esc"):
        running = False
        break
\end{lstlisting}

\textbf{Our control scheme (similar to video game controls):}
\begin{itemize}
	\item \texttt{W/S}: Move forward/backward
	\item \texttt{A/D}: Strafe left/right
	\item \texttt{Z/C}: Go up/down
	\item \texttt{Q/E}: Rotate left/right
	\item \texttt{ESC}: Exit the program
\end{itemize}

The cool thing about body-frame control is that "forward" always means the direction the drone is facing. So if we turn 90 degrees and press W, the drone moves in its new forward direction—just like in a first-person game.

\textbf{6. Sending Commands to the Drone:}
\begin{lstlisting}[language=Python]
    # Make the drone rotate
    if YawRate != 0:
        client.rotateByYawRateAsync(YawRate, 0.1)
        
    # Make the drone move
    if vx != 0 or vy != 0 or vz != 0:
        client.moveByVelocityBodyFrameAsync(vx, vy, vz, 0.1)
    elif YawRate == 0:
        # Hold position when we're not pressing anything
        client.moveByVelocityBodyFrameAsync(0, 0, 0, 0.1)
\end{lstlisting}

Each command lasts 0.1 seconds, which creates smooth, continuous motion. When we're not pressing any keys, we send a "zero velocity" command to keep the drone from drifting.

\textbf{7. Recording Flight Data:}
\begin{lstlisting}[language=Python]
    state = client.getMultirotorState()
    position = state.kinematics_estimated.position
    orientation = state.kinematics_estimated.orientation
    gps_data = client.getGpsData()
    
    roll, pitch, yaw = airsim.to_eularian_angles(orientation)
    
    # Save everything to our log file
    current_time = datetime.now().strftime('%H:%M:%S.%f')[:-3]
    writer.writerow([
        current_time, currentCmd, vx, vy, vz, YawRate,
        position.x_val, position.y_val, position.z_val,
        math.degrees(roll), math.degrees(pitch), math.degrees(yaw),
        gps_data.gnss.geo_point.latitude,
        gps_data.gnss.geo_point.longitude,
        gps_data.gnss.geo_point.altitude
    ])
    
    logfile.flush()
\end{lstlisting}

\textbf{What we're recording:}
\begin{itemize}
	\item The drone's position in 3D space (where it is)
	\item Its orientation (which way it's tilted and facing)
	\item GPS coordinates (latitude, longitude, altitude)
	\item What command we just gave it
	\item The exact time
\end{itemize}

The \texttt{flush()} command ensures our data gets written to the file immediately, so we don't lose anything if the program crashes.

\textbf{8. Cleanup:}
\begin{lstlisting}[language=Python]
logfile.close()
print("Flight log saved to:", log_path)
\end{lstlisting}

When we exit the program, we close the log file and print where it was saved.

\subsubsection{Visualizing the Flight Data}

After flying the drone, we wanted to see what happened during the flight. We wrote a second script to create 14 different graphs showing various aspects of the flight.

\textbf{Key features of our visualization tool:}

\textbf{1. Finding and Loading Flight Logs:}
\begin{lstlisting}[language=Python]
def get_latest_log_file():
    """Find the most recent flight recording"""
    log_files = glob.glob("flightLogs/drone_path_*.csv")
    if not log_files:
        return None
    return max(log_files, key=os.path.getctime)

def list_available_logs():
    """Show all available flight recordings"""
    log_files = glob.glob("flightLogs/drone_path_*.csv")
    # ... displays timestamps of each flight ...
\end{lstlisting}

The script can automatically find our most recent flight, or we can choose to analyze a specific earlier flight.

\textbf{2. Making the Data Easier to Understand:}
\begin{lstlisting}[language=Python]
# Convert to more intuitive altitude values
df['z_drone'] = -df['z']   # Now positive means higher altitude
df['vz_drone'] = -df['vz'] # Positive velocity means going up
\end{lstlisting}

Since the NED coordinate system can be confusing, we flip the Z values so that positive means "up"—which is how most people naturally think about altitude.

\textbf{3. The 14 Graphs We Generate:}

Our script creates a comprehensive set of plots:

\begin{enumerate}
	\item \textbf{X Position Over Time}: Shows how far north or south we flew
	\item \textbf{Y Position Over Time}: Shows how far east or west we flew
	\item \textbf{Z Position (Altitude) Over Time}: Shows our height changes
	\item \textbf{Velocity X Over Time}: Our forward/backward speed
	\item \textbf{Velocity Y Over Time}: Our left/right speed
	\item \textbf{Velocity Z Over Time}: Our up/down speed
	\item \textbf{Yaw Rate Over Time}: How fast we were spinning
	\item \textbf{Roll Angle Over Time}: How much we were banking left or right
	\item \textbf{Pitch Angle Over Time}: How much the nose was up or down
	\item \textbf{Yaw Angle Over Time}: Which direction we were facing
	\item \textbf{GPS Latitude Over Time}: Our north-south geographic position
	\item \textbf{GPS Longitude Over Time}: Our east-west geographic position
	\item \textbf{GPS Altitude Over Time}: Our altitude above sea level
	\item \textbf{Flight Trajectory (Top View)}: A bird's-eye view of our entire path
\end{enumerate}

\textbf{4. Creating an Informative Trajectory Map:}
\begin{lstlisting}[language=Python]
# Color the path based on time (darker = later in flight)
scatter = plt.scatter(df['x'], df['y'], c=time_seconds, cmap='viridis', 
                     s=30, alpha=0.6, edgecolors='black', linewidth=0.5)

# Draw the flight path as a line
plt.plot(df['x'], df['y'], 'b-', alpha=0.3, linewidth=1.5, label='Flight Path')

# Mark where we started (green circle) and ended (red X)
plt.scatter(df['x'].iloc[0], df['y'].iloc[0], color='green', s=300, 
           marker='o', edgecolors='black', linewidth=2, label='Start', zorder=5)
plt.scatter(df['x'].iloc[-1], df['y'].iloc[-1], color='red', s=300, 
           marker='X', edgecolors='black', linewidth=2, label='End', zorder=5)

# Add arrows to show which direction we were flying
arrow_interval = len(df) // 10
for i in range(0, len(df) - arrow_interval, arrow_interval):
    dx = df['x'].iloc[i + arrow_interval] - df['x'].iloc[i]
    dy = df['y'].iloc[i + arrow_interval] - df['y'].iloc[i]
    if abs(dx) > 0.01 or abs(dy) > 0.01:
        plt.arrow(df['x'].iloc[i], df['y'].iloc[i], dx, dy, 
                 head_width=0.5, head_length=0.3, fc='blue', ec='blue', 
                 alpha=0.4, length_includes_head=True)
\end{lstlisting}

This creates a really nice visualization where:
\begin{itemize}
	\item The color gradient shows time progression (purple → yellow as time goes on)
	\item A green circle marks our starting point
	\item A red X marks where we ended
	\item Blue arrows point in the direction of flight
	\item The scale is accurate so we can see real distances
\end{itemize}

\textbf{5. Calculating Flight Statistics:}
\begin{lstlisting}[language=Python]
speed = np.sqrt(df['vx']**2 + df['vy']**2 + df['vz']**2)
dx = np.diff(df['x'], prepend=df['x'].iloc[0])
dy = np.diff(df['y'], prepend=df['y'].iloc[0])
dz = np.diff(df['z'], prepend=df['z'].iloc[0])
distances = np.sqrt(dx**2 + dy**2 + dz**2)
total_distance = np.sum(distances)

print(f"Max speed: {speed.max():.2f} m/s")
print(f"Average speed: {speed.mean():.2f} m/s")
print(f"Total distance traveled: {total_distance:.2f} m")
print(f"Max altitude change: {df['z_drone'].max() - df['z_drone'].min():.2f} m")
\end{lstlisting}

The script also calculates some useful statistics:
\begin{itemize}
	\item Our top speed during the flight
	\item Our average speed
	\item Total distance we covered (accounting for all three dimensions)
	\item How much our altitude varied
\end{itemize}

\subsubsection{How to Run Everything}

\textbf{Step 1: Launch AirSim}
\begin{itemize}
	\item Start the Unreal Engine environment with the AirSim plugin
	\item Make sure the drone appears at the starting position we configured
	\item Check that everything is running smoothly
\end{itemize}

\textbf{Step 2: Fly the Drone}
\begin{verbatim}
python drone_control.py
\end{verbatim}

\begin{itemize}
	\item Press \texttt{Y} to arm the motors and take off
	\item Use the keyboard controls to fly around
	\item Press \texttt{L} when you're ready to land
	\item Press \texttt{ESC} to exit and save your flight log
\end{itemize}

\textbf{Step 3: Analyze Your Flight}
\begin{verbatim}
python plot_flight_data.py
\end{verbatim}

\begin{itemize}
	\item Choose option 1 to analyze your most recent flight
	\item Or choose option 2 to pick a specific flight log
	\item Close each graph window to see the next one
	\item All 14 graphs will display one after another
\end{itemize}

\subsubsection{What the Graphs Show Us}

The plots give us a complete picture of how the drone performed:

\begin{figure}[h]
	\centering
	\begin{minipage}[t]{0.48\textwidth}
		\centering
		\includegraphics[width=\textwidth]{Figure_1.png}
		\caption*{(a) X Position Over Time}
	\end{minipage}\hfill
	\begin{minipage}[t]{0.48\textwidth}
		\centering
		\includegraphics[width=\textwidth]{Figure_2.png}
		\caption*{(b) Y Position Over Time}
	\end{minipage}
	\caption{Position data showing how far we traveled in each direction}
	\label{fig:lab5-position}
\end{figure}

\begin{figure}[h]
	\centering
	\begin{minipage}[t]{0.48\textwidth}
		\centering
		\includegraphics[width=\textwidth]{Figure_3.png}
		\caption*{(c) Z Position (Altitude) Over Time}
	\end{minipage}\hfill
	\begin{minipage}[t]{0.48\textwidth}
		\centering
		\includegraphics[width=\textwidth]{Figure_14.png}
		\caption*{(d) Flight Trajectory (Top View)}
	\end{minipage}
	\caption{Altitude changes and the complete flight path from above}
	\label{fig:lab5-altitude-trajectory}
\end{figure}

\begin{figure}[h]
	\centering
	\begin{minipage}[t]{0.48\textwidth}
		\centering
		\includegraphics[width=\textwidth]{Figure_4.png}
		\caption*{(e) Velocity X Over Time}
	\end{minipage}\hfill
	\begin{minipage}[t]{0.48\textwidth}
		\centering
		\includegraphics[width=\textwidth]{Figure_5.png}
		\caption*{(f) Velocity Y Over Time}
	\end{minipage}
	\caption{Speed variations showing how the drone responded to our controls}
	\label{fig:lab5-velocity}
\end{figure}

\begin{figure}[h]
	\centering
	\begin{minipage}[t]{0.48\textwidth}
		\centering
		\includegraphics[width=\textwidth]{Figure_6.png}
		\caption*{(g) Velocity Z Over Time}
	\end{minipage}\hfill
	\begin{minipage}[t]{0.48\textwidth}
		\centering
		\includegraphics[width=\textwidth]{Figure_7.png}
		\caption*{(h) Yaw Rate Over Time}
	\end{minipage}
	\caption{Vertical speed and rotation rate showing climb/descent and turning maneuvers}
	\label{fig:lab5-vz-yawrate}
\end{figure}

\begin{figure}[h]
	\centering
	\begin{minipage}[t]{0.48\textwidth}
		\centering
		\includegraphics[width=\textwidth]{Figure_8.png}
		\caption*{(i) Roll Angle Over Time}
	\end{minipage}\hfill
	\begin{minipage}[t]{0.48\textwidth}
		\centering
		\includegraphics[width=\textwidth]{Figure_9.png}
		\caption*{(j) Pitch Angle Over Time}
	\end{minipage}
	\caption{Roll and pitch angles showing how the drone tilted during maneuvers}
	\label{fig:lab5-roll-pitch}
\end{figure}

\begin{figure}[h]
	\centering
	\begin{minipage}[t]{0.48\textwidth}
		\centering
		\includegraphics[width=\textwidth]{Figure_10.png}
		\caption*{(k) Yaw Angle Over Time}
	\end{minipage}\hfill
	\begin{minipage}[t]{0.48\textwidth}
		\centering
		\includegraphics[width=\textwidth]{Figure_11.png}
		\caption*{(l) GPS Latitude Over Time}
	\end{minipage}
	\caption{Heading direction and geographic north-south position changes}
	\label{fig:lab5-yaw-lat}
\end{figure}

\begin{figure}[h]
	\centering
	\begin{minipage}[t]{0.48\textwidth}
		\centering
		\includegraphics[width=\textwidth]{Figure_12.png}
		\caption*{(m) GPS Longitude Over Time}
	\end{minipage}\hfill
	\begin{minipage}[t]{0.48\textwidth}
		\centering
		\includegraphics[width=\textwidth]{Figure_13.png}
		\caption*{(n) GPS Altitude Over Time}
	\end{minipage}
	\caption{Geographic east-west position and altitude above sea level from GPS}
	\label{fig:lab5-lon-alt}
\end{figure}

\clearpage

\subsection{What We Accomplished}
\begin{itemize}
	\item Successfully set up a realistic drone simulation environment
	\item Configured a virtual quadcopter with a front-facing camera
	\item Positioned the drone at a sensible starting point just above the ground
	\item Added realistic wind conditions to make the simulation more challenging
	\item Built a keyboard control system that lets us fly in 8 directions plus rotation
	\item Achieved smooth flight at 2.5 m/s with responsive 25 deg/s turning
	\item Created a comprehensive logging system that records 15 different parameters 20 times per second
	\item Captured position, velocity, and orientation data throughout each flight
	\item Recorded GPS coordinates for geographic reference
	\item Logged all our control inputs with precise timestamps
	\item Successfully flew the drone from takeoff to landing at various target locations
	\item Generated 14 detailed graphs showing every aspect of flight performance
	\item Created a bird's-eye trajectory map of the complete flight path
	\item Calculated useful statistics like maximum speed, average speed, and total distance traveled
	\item Analyzed altitude variations and GPS coordinate changes
	\item Visualized how the drone's state evolved over time in response to our inputs
\end{itemize}

\subsection{What We Learned}
\begin{itemize}
	\item AirSim uses the NED coordinate system where Z is inverted—it takes some getting used to, but it's a standard in aviation
	\item The \texttt{settings.json} file is crucial and must be in exactly the right location for AirSim to find it
	\item The SimpleFlight controller is great for learning because it provides stable, predictable flight
	\item Auto-creation is convenient—the drone appears automatically without manual placement
	\item Wind simulation adds realism and teaches us how drones must constantly adjust to stay on course
	\item We need to explicitly enable API control before our code can command the drone
	\item Just like real drones, we must arm the motors before takeoff for safety
	\item The automatic takeoff function handles the tricky part of getting airborne safely
	\item Body-frame velocity control feels natural, like a first-person game—forward is always relative to where the drone is facing
	\item Commands in body frame move relative to the drone's orientation, not the world
	\item Short command durations (0.1 seconds) at high frequency create smooth, responsive control
	\item We can rotate and move at the same time, enabling complex maneuvers
	\item Sending zero velocity when idle is important—otherwise wind and momentum cause drift
	\item The state data includes everything we need to know about where the drone is and how it's moving
	\item Orientation comes as a quaternion, which we need to convert to roll, pitch, and yaw angles to make sense of it
	\item GPS coordinates are separate from the local XYZ position—they give us a global reference
	\item Recording data 20 times per second captures plenty of detail without overwhelming us with information
	\item The flush command prevents data loss—without it, we might lose recent data if something goes wrong
	\item Time-series plots reveal cause and effect—we can see exactly how our inputs affected the drone
	\item Velocity graphs show us how quickly the drone responds and whether control is smooth
	\item Orientation plots tell us about stability—wild swings would indicate problems
	\item The trajectory plot with time-based coloring makes it easy to follow the sequence of events
	\item Direction arrows clarify which way we were flying at different points
	\item Calculating total distance requires 3D math—we can't just look at X and Y
	\item Flipping the Z values in our plots makes altitude easier to understand for humans
	\item GPS gives us real-world coordinates that complement the local position data
	\item Having multiple control modes (velocity, position, angle) lets us choose the right tool for the job
	\item Fast keyboard polling (checking 20 times per second) makes control feel responsive
	\item Separate takeoff and landing procedures ensure we don't accidentally crash during startup or shutdown
	\item Good visualization is essential—raw numbers don't tell the story nearly as well as graphs do
\end{itemize}

\subsection{Conclusions}

This lab gave us hands-on experience with drone simulation and control. We successfully set up AirSim with realistic physics, wind, and sensors. The configuration file let us customize everything from the drone's starting position to camera placement.

Our Python control script turned our keyboard into a drone controller, using intuitive body-frame commands that made flying feel natural. We could smoothly navigate from any starting point to any destination using the WASD+QEZC control scheme.

The logging system we built captured every detail of our flights at 20 samples per second. This gave us a complete record of positions, velocities, orientations, and GPS coordinates. When we analyzed this data with our visualization script, we got 14 different views into how the drone behaved, plus a comprehensive trajectory map.

The trajectory plot was particularly useful—it showed not just where we went, but when and in which direction. The color gradient and arrows made it easy to understand the sequence of our flight.

What stands out from this experience is how important it is to understand coordinate systems (especially that tricky inverted Z-axis!), how body-frame control makes piloting intuitive, and how valuable detailed logging is for understanding what actually happened during a flight. The combination of real-time flying and after-flight analysis gave us a complete picture of drone navigation.

This lab provided a solid foundation for more advanced drone work. We can now experiment with autonomous navigation, test control algorithms, and develop more sophisticated flight behaviors—all in a safe simulated environment where crashes don't cost us thousands of dollars!

\clearpage